\documentclass[10pt,a4paper]{amsart}
\usepackage[margin=19mm]{geometry}
\usepackage{amsmath,amssymb,mathtools}
\usepackage[scr=rsfs,cal=euler]{mathalfa}
\usepackage{xfrac}
\usepackage[unicode,hidelinks,bookmarks=false]{hyperref}
\newcommand{\ie}{i.e.\ }
\newcommand{\cf}{cf.\ }
\newcommand{\resp}{respectively\ }
\newcommand{\Subsection}{\subsection}
\providecommand{\nobreakdash}{\nobreak\hbox{-}}
\setlength{\emergencystretch}{2em}
\multlinegap0pt
\usepackage{mathrsfs,latexsym,amscd,tikz-cd}
\usepackage{fullpage}
\newcommand{\D}{\mathbb{D}}
\newcommand{\N}{\mathbb{N}}
\newcommand{\Z}{\mathbb{Z}}
\newcommand{\Q}{\mathbb{Q}}
\newcommand{\R}{\mathbb{R}}
\newcommand{\p}{\mathbb{P}}
\newcommand{\A}{\mathbb{A}}
\newcommand{\C}{\mathbb{C}}
\newcommand{\Int}{\operatorname{Int}}
\newtheorem{thm}{Theorem}[section]
\newtheorem{cor}{Corollary}[section]
\newtheorem{exercise}[thm]{Exercise}
\newtheorem{ax}{Axiom}
\newtheorem*{defn}{Definition}
\newtheorem*{solution}{Solution}
\newtheorem*{example}{Example}
\newtheorem*{examples}{Examples}
\setlength{\parskip}{10pt}
\begin{document}
\title[Combinatorial Hodge Index Theorem for Polytopes]{Combinatorial Hodge Index Theorem for Polytopes}
\author{Jacob B. Wood}
\date{}
\maketitle
\noindent\emph{Source and licence.} Combinatorial Hodge Index Theorem for Polytopes. Version arXiv:2608.18003v1 (CC BY 4.0). This material is distributed under the Creative Commons Attribution 4.0 International licence, http://creativecommons.org/licenses/by/4.0/, as recorded in the arXiv OAI licence metadata for this exact version and shown on the abstract page https://arxiv.org/abs/2608.18003v1. Full rights notice: licenses/arxiv-2608.18003-CC-BY-4.0.txt.\par\smallskip
\noindent\textit{Editorial scope.} Counted unit: the original \texttt{thm} environment \texttt{th33} at source lines 133--135 together with its complete proof at lines 138--184; the delivered document keeps source lines 113--185 so that every referenced label and every citation resolves inside it. Native editable source is the paper's own concatenated TeX; the AMS wrapper is editorial formatting only. No publisher class, upstream script or unrelated artwork is redistributed.
\section{The Combinatorial K\"ahler Package and signature of fans}

For toric varieties, there is a Lefschetz operator on intersection cohomology, and we would like to find such an operator for combinatorial intersection cohomology. It is unknown whether there is such an operator for general fans. However, when $\Phi$ admits a strictly convex piecewise linear function (e.g., the inner normal fan of a convex polytope $P$), then that function as an element of $\mathcal{A}_{\Phi}(\Phi)$ is a Lefschetz operator for $IH(\Phi)$. Fans which have a strictly convex piecewise linear function are called \textbf{projective} because when the fan has an associated toric variety, this condition is equivalent to the toric variety being projective (see \cite[6.9]{D}). When $\Phi$ is projective, Karu \cite[Theorem 0.1]{K} showed:


\begin{thm} (Hard Lefschetz) \label{th31} Let $\Phi$ be a projective fan with a strictly convex piecewise linear function $\ell$. Then the map \[\ell^k:IH^{n-k}(\Phi)\rightarrow IH^{n+k}(\Phi)\] is an isomorphism for all $k\geq 1.$
\end{thm}

Observe the following consequences of the hard Lefschetz theorem. First, the (intersection) Betti numbers of a fan $\Phi$, which we will denote \[ih^j_\Phi:=\dim IH^j(\Phi),\] form a $\mathbf{symmetric}$ sequence, so \[ih_\Phi^j=ih_\Phi^{2n-j}.\] Furthermore, since the composition of functions $\ell^{k}:IH^{n-k}(\Phi)\rightarrow IH^{n+k}(\Phi)$ is an isomorphism, the first iteration $\ell:IH^{n-j}(\Phi)\rightarrow IH^{n-j+2}(\Phi)$ is an injection, and the last map $\ell:IH^{n+j-2}(\Phi)\rightarrow IH^{n+j}(\Phi)$ is a surjection. Thus, the sequence of even index Betti numbers $ih_\Phi^{2j}$ is $\mathbf{unimodal}$. The kernel of the surjection $\ell^{j+1}:IH^{n-j}(\Phi)\rightarrow IH^{n+j+2}(\Phi)$ is called the $\mathbf{primitive\ subspace}$ of $IH^{n-j}(\Phi)$, which we denote \[\text{Prim}_\ell IH^{n-j}(\Phi)=\text{ker}(\ell^{j+1}|_{IH^{n-j}(\Phi)}).\] For the dimensions of the primitive subspaces, we write \[ip_\Phi^{n-j}:=\dim \text{Prim}_\ell IH^{n-j}(\Phi)=ih_\Phi^{n-j}-ih_\Phi^{n-j-2}.\] It is sometimes convenient to encode these invariants as the coefficients of a polynomial, so we define: $$ih_\Phi(q):=\sum_{j=0}^{2n} ih_\Phi^j q^j$$
$$ip_\Phi(q):=\sum_{j=0}^{n} ip_\Phi^j q^j.$$


The Hodge-Riemann relations for fans proven by Bressler-Lunts \cite{BL2} describe how the cohomology pairing, $(\cdot,\cdot)_\Phi:IH^{n-k}(\Phi)\times IH^{n+k}(\Phi)\rightarrow IH^{2n}(\Phi)$, behaves on each primitive summand.


\begin{thm} (Hodge-Riemann Relations) Let $\Phi$ and $\ell$ be as in Theorem \ref{th31}. Then for any non-zero $a\in \text{Prim}_\ell IH^{n-k}(\Phi)$ we have $$Q_\Phi(a):=(-1)^\frac{n-k}{2}(a,\ell^ka)_\Phi>0.$$
\end{thm}

Combining the above two theorems, we obtain the following result analogous to the proof of the classical Hodge index theorem \cite{H}. This gives a combinatorial proof of a similar formula obtained by Maxim-Sch\"urmann \cite[Example 5.8, Example 6.7]{MS} in the toric context (i.e., for rational polytopes). Note that in the notations of \cite{MS}, our result is stated for their polar polytope, see \cite[Formula (6.13)]{MS}, which explains the difference in powers of $(-2)$ in our formulae. 


\begin{thm} (Signature of a projective fan) \label{th33} Let $\Phi$ be the projective fan of dimension $n$ associated to a convex polytope. The signature of its intersection pairing is given by $$\sigma(\Phi)=\sum_{\Delta\prec\Phi} g_\Delta(-1)(-2)^{n-\dim\Delta-1}$$
where $g_\Delta(t)$ is Stanley's $g$-polynomial associated to the proper face $\Delta$.
\end{thm}

\begin{proof} We will begin by computing the signature $\sigma(\Phi):=\sigma(IH^n(\Phi))$ in terms of the Betti numbers of $\Phi$. Then we use the fact that the Betti numbers coincide with certain combinatorial invariants.

When $n$ is odd, the middle intersection cohomology group is $IH^n(\Phi)=0$ because all odd dimensional intersection cohomology vanishes, so $\sigma(\Phi)=0$. Observe that for odd $n$, hard Lefschetz (Theorem \ref{th31}) implies symmetry of the non-zero Betti numbers: $ih^{2k}=ih^{2n-2k}$. Thus, 
\begin{equation}\sum_{k=0}^n (-1)^kih_{\Phi}^{2k}=0=\sigma(\Phi).
\end{equation}

When $n$ is even, we obtain the following orthogonal decomposition of $IH^n(\Phi)$ by the hard Leschetz theorem:
\begin{equation}
   IH^n(\Phi)=\bigoplus_k \ell^{k}\text{Prim}_\ell IH^{n-2k}(\Phi), 
\end{equation}
so we have
\begin{equation}
\begin{split}
    \sigma(\Phi)&=\sum_{0\leq k\leq \frac{n}{2}}\sigma(\text{Prim}_\ell IH^{n-2k}(\Phi)) \\
   & \overset{(HR)}{=}\sum_{0\leq k\leq \frac{n}{2}}(-1)^\frac{n-2k}{2}\text{dim}(\text{Prim}_\ell IH^{n-2k}(\Phi)),\\
&=\sum_{0\leq k\leq \frac{n}{2}}(-1)^\frac{n-2k}{2}(ih_\Phi^{n-2k}-ih_\Phi^{n-2k-2}),
\end{split}
\end{equation}
where the sign of the terms in the second equality is given by the Hodge-Riemann relations. Furthermore, rearranging terms shows 

\begin{equation}\label{eq1}
    \sum_{0\leq k\leq \frac{n}{2}}(-1)^\frac{n-2k}{2}(ih_\Phi^{n-2k}-ih_\Phi^{n-2k-2})=\sum_{0\leq k\leq \frac{n}{2}}(-1)^\frac{n-2k}{2}ih_\Phi^{n-2k}+\sum_{0\leq k\leq \frac{n}{2}}(-1)^{\frac{n-2k}{2}+1}ih_\Phi^{n-2k-2}
\end{equation}
Then reindexing the second summand of the right side of Equation \ref{eq1}, we have that
\begin{equation*}
\begin{split}
    \sum_{0\leq k\leq \frac{n}{2}}(-1)^\frac{n-2k}{2}ih_\Phi^{n-2k}&+\sum_{0\leq k\leq \frac{n}{2}}(-1)^{\frac{n-2k}{2}+1}ih_\Phi^{n-2k-2}=2\left(\sum_{0\leq k< \frac{n}{2}}(-1)^\frac{n-2k}{2}ih_\Phi^{n-2k}\right)+(-1)^{\frac{n}{2}} ih_\Phi^n \\
&   \overset{(HL)}{=}\sum_{0\leq k< \frac{n}{2}}(-1)^\frac{n-2k}{2}ih_\Phi^{n-2k}+\sum_{0\leq k< \frac{n}{2}}(-1)^\frac{n+2k}{2}ih_\Phi^{n+2k}+(-1)^{\frac{n}{2}} ih_\Phi^n \\
&    =\sum_{0\leq r\leq n}(-1)^rih_\Phi^{2r},
\end{split}
\end{equation*}
where the equality labeled (HL) follows from the hard Lefschetz theorem, which implies the symmetry of Betti numbers.

From now on, we treat both the even and the odd cases together. Bressler-Lunts show that the Hard Lefschetz theorem implies that local $IH$ Betti numbers are well-defined \cite[Corollary 7.4]{BL} and that the following formulae \cite[Proposition 7.7]{BL} relate $ih_\Phi(q)$ to $ip_\Phi(q)$:
\[ip_o(q)=ih_o(q)\equiv1\]
\[ip^j_\sigma=\begin{cases}
                ih_\sigma^j-ih_\sigma^{j-2} \text{ for } 0\leq j\leq d  \\
                0 \text{ otherwise}
                 \end{cases}
\]
\[ih_\sigma(q)=\sum_{\Delta\in\sigma}(q^2-1)^{d-\text{dim}\Delta-1}ip_\Delta(q).\]
They then compare with Stanley's $h$- and $g$-polynomials whose relations coincide with those of the $ip$ and $ih$ polynomials of $\Phi$ when we declare the poset to be the poset of cones under inclusion and the poset's rank function to be the dimension of the cone minus 1. Thus, the pairs of polynomials are defined by the same inductive procedure with the same base case, so they must have the same coefficients, and $$h_\Phi(q^2)=ih_\Phi(q),\ \text{{and}}$$
$$g_\Phi(q^2)=ip_\Phi(q).$$
Then the signature $\sigma(\Phi)$ can be rewritten as
$$\sigma(\Phi)=\sum_i(-1)^iih_\Phi^{2i}=\sum_i(-1)^ih_\Phi^{i}=\sum_{\Delta\in\Phi} g_\Delta(-1)(-2)^{d-\text{dim}\Delta-1},$$
which completes the proof.
\end{proof}


\begin{thebibliography}{99}
\bibitem[BL]{BL} P. Bressler, V. A. Lunts, {\it Intersection Cohomology on Nonrational Polytopes}, Compositio Mathematica 135 (2003), 245-278
\bibitem[BL2]{BL2} P. Bressler, V. Lunts, {\it Hard Lefschetz theorem and Hodge-Riemann relations for intersection cohomology of nonrational polytopes}, Indiana Univ. Math. J., 54(1) (2005), 263–307
\bibitem[D]{D} V. I. Danilov, {\it The Geometry of Toric Varieties}, Russ. Math. Surv. 33 97 (1978)
\bibitem[H]{H} F. Hirzebruch, {\it Topological methods in algebraic geometry}, Springer, 1966
\bibitem[K]{K} K. Karu, {\it Hard Lefschetz theorem for nonrational polytopes}, Invent. Math. 157 (2004), 419–447
\bibitem[MS]{MS} L. Maxim, J. Sch\"urmann, {\it Weighted Ehrhart theory via mixed Hodge modules on toric varieties}, International Mathematics Research Notices, Vol. 2025, Issue 7, pages 1-25
\end{thebibliography}
\end{document}
